\section{Post-saturation unresolved-route authority}
\label{sec:p2x}

The post-saturation literature search raises the comparison standard rather than shrinking the systems question. Verification-aware claim coverage, decision-sufficiency and utility stopping, coverage-driven deep-research completion, finite source-obligation certificates, and generic workflow-obligation graphs are all granted as donor mechanisms \cite{halt2026,knowplan2026,donotstopearly2026,confidencebasedstop2026,icore2026}. P2-X then asks the remaining scientific question: after those mechanisms have extracted everything they can from the routes that are actually available, what authority exists to close the task when a \emph{material} route is unavailable, censored, provider-invalid, or still carries an unresolved question-conditioned processing obligation?

\paragraph{Exact contract battery.}
P2-X freezes five acquisition families---scientific literature, web/API evidence, dataset registries, code/artifact repositories, and experiment/measurement routes---with eight hostile archetypes per family. The cases distinguish all-obligations-discharged closure, one available route remaining open, unavailable/censored/provider-invalid material routes, duplicate route coverage, already-acquired but not question-processed content, and low-utility hard routes that remain open. Route-local decisions are \texttt{ROUTE\_STOP}, \texttt{ROUTE\_CONTINUE}, or \texttt{ROUTE\_CANNOT\_CHECK}; task terminals are \texttt{TASK\_STOP}, \texttt{CONTINUE}, or \texttt{CANNOT\_CHECK}. Route order is seed-permuted, and held-out variants add unavailable \emph{non-material} decoy routes so that blanket ``any unavailable source blocks'' logic cannot pass.

\paragraph{Strong donor products.}
The primary B1 comparator is deliberately donor-complete. It receives route identity/materiality, HALT-like evidence coverage, decision sufficiency, utility stopping, finite route certificates, deduplication, route state, and the question-conditioned processing coordinate. It also makes the same correct route-local decisions as P2-X. Its only systematic difference is global aggregation: once all \emph{available} material routes locally stop and their acquired content is processed, it treats the task as closed, excluding unavailable/censored/provider-invalid material routes from the closure denominator. B2 allows a global sufficiency signal to authorize task stop. B3 is an ideal information-equivalent product given the exact P2-X unresolved-route semantics and is expected to tie; it tests portability of the authority rule rather than serving as a weak baseline.

\paragraph{Protected result: exact closure advantage.}
On 400 prospectively frozen protected cases, P2-X obtains \textbf{400/400 exact scientific closure decisions} (ESCD $=1.000$), versus 250/400 ($0.625$) for the donor-complete B1 product and 100/400 ($0.250$) for B2. The paired P2-X--B1 advantage is \textbf{$+0.375$} with a predeclared domain-stratified bootstrap 95\% interval $[0.3275,0.4225]$, above the registered $+0.10$ practical margin; exact McNemar discordance is $b=150,c=0$. P2-X records \textbf{zero false task closures}, preserves 1.000 clean task-stop accuracy when every material obligation is discharged, and has a positive advantage over B1 in all five acquisition families. B1's 150 ESCD errors are exactly the three unresolved-material-route archetypes---unavailable, censored, and provider-invalid---while it passes every other archetype.

An independent implementation reproduces all 400 case identities, protected-bundle and canonical decision-row hashes, every headline arm count, and zero B3 decision mismatches. The non-material unavailable decoys cause no P2-X materiality error, and all statistics include the held-out decoy variants.

\paragraph{Portability result.}
B3 also obtains 400/400 and is decision-identical to P2-X. This is a positive implementation result: the unresolved-route authority rule is sufficient and portable. It can be realized as the ORION closure layer or as a correctly enriched donor product without changing any registered decision. The P2-X advantage over B1 therefore isolates a specific missing coupling in a realistic donor-complete interface rather than depending on centralized expressivity or on weak stopping machinery.

The historical OpenAIRE/Crossref Wide campaign remains a separate provider-validity record and is not used to support the P2-X effect. The present claim is the protected exact-contract result: when a material acquisition route remains unresolved, route-local stopping is not task-global scientific closure, and making that authority distinction explicit yields perfect registered closure decisions against the strong donor-complete comparator.
